\paragraph{Opposing changes can cancel.}
The CALM flip rates show that duplication frequently changes which
answer the judge prefers. JudgeBench reveals the consequences of
those changes. In the primary comparison, duplicating the incorrect
answer introduced 94 wrong judgments and one tie among previously
correct verdicts, and recovered 88 previously incorrect verdicts.
These opposing transitions explain why the aggregate accuracy moved
only half a percentage point. For a system that uses the judge to make
individual choices, both directions matter.

\paragraph{The rubric changes which side benefits.}
With ordered scoring, the 4B rubric reduces accuracy by 9.0 points
when the correct answer is duplicated and raises it by 4.5 points
when the incorrect answer is duplicated. Duplicating both answers
changes accuracy by less than one point in that setting. Pooling
amplifies the one-sided effects under the rubric. Its accuracy rises
from 55.0\% to 69.5\% when the incorrect answer is duplicated and
falls to 40.5\% when the correct answer is duplicated. These results
are consistent with a penalty for the repeated answer under this
prompt. The rubric can therefore improve the score for one style
placement while degrading it for the mirror placement.

\paragraph{Higher accuracy can coexist with instability.}
Pooling the two answer orders under the brief prompt raised 4B
accuracy on original JudgeBench pairs from 53.7\% to 58.6\%. It
also raised accuracy for incorrect-answer duplication from 53.1\%
to 56.4\%. The resulting loss from duplication was 2.2 points for
the pooled rule, compared with 0.5 for the ordered rule, and the
estimated difference had an interval spanning zero. Similarly,
forcing 0.6B to select an answer raises its coverage to 100\% on
original labeled pairs. Its pooled forced-choice accuracy remains
close to 50\%. Accuracy, conditional error transitions, and coverage
answer distinct questions about each protocol.

\paragraph{Precision affects abstention.}
On the 90 matched JudgeBench pairs, switching the 0.6B checkpoint from
NF4 to BF16 raises brief-prompt coverage on original comparisons from
47.2\% to 80.0\%. Overall accuracy rises from 25.0\% to 41.1\%.
Accuracy among selected answers is 52.9\% with NF4 and 51.4\%
with BF16. Much of the overall accuracy change on this subset
comes from the smaller share of abstentions.

\paragraph{Concrete errors.}
In JudgeBench pair \texttt{01c32337}, the incorrect candidate on an
ABO allele-frequency question gives computed values that disagree
with its selected choice. Duplicating that candidate causes three
correct-to-wrong transitions across the 1.7B and 4B configurations.
In pair \texttt{ff5f3ce7}, the benchmark credits an answer about
tribute in the Aztec Empire. Adding \emph{Certainly.} to a competing
explanation based on religious sacrifice produces three
correct-to-wrong transitions. These cases retain their original
content under the edits.

\paragraph{Limitations.}
The benchmark labels evaluate correctness on the included questions.
The same texts could receive different rankings under a rubric
that explicitly values concision or ease of reading. A duplicated
answer occupies more positions in the prompt, so the repetition
intervention changes both redundancy and token placement. The
confidence and formatting treatments each use one fixed edit.
Some questions still specify an answer format. For example, a
reasoning question about cuts through a pentagon requests a bold
integer, and an edit also introduces bold text elsewhere. The
exploratory screen checks how these cases affect the aggregate
comparison. The three judges belong to one model family, and
all reported effects are conditional on the selected checkpoint,
prompt, and decoding rule.
